\section{Discussion}
\subsection{What the results establish}
The frozen source selector improves the paired-seed waiting outcome over MaxPressure in the two-network comparison. Its Jinan improvement over the legacy target anchor and exact fallback to that anchor in Los Angeles support the tested source-selection procedure. The favorable PhasePressure point estimate gives a less conclusive comparison because its interval crosses zero.

The main comparison and the development tests in Appendix~\ref{app:development-results} support the tested predictor configurations over the dense linear residual baseline. They do not show that every component is necessary. In particular, the gain may reflect the nonlinear regressor, action normalization, feature restrictions, pairwise learning, weighting, or their interaction. A capacity-matched nonlinear baseline and controlled component ablations are needed before attributing the effect to a specific architectural choice. The reference-only comparison provides one such partial ablation, but its closed-loop interval remains inconclusive.

\subsection{Why offline ranking may not predict control performance}
There are two changes between the supervised task and deployment. First, the labels measure one focal phase choice followed by PhasePressure, with other intersections also using PhasePressure. The deployed controller repeatedly chooses learned phases at multiple intersections. It can therefore visit states and joint action combinations that differ from those under the label continuation. Second, the training target is a 60-s vehicle-load measure, whereas evaluation measures waiting and completion over an hour. An action can improve the former without improving the latter. These differences offer plausible explanations for the imperfect relationship between offline regret and closed-loop waiting, but the present experiments do not identify their individual causal contributions.

\subsection{Scope of the evidence}
This result includes eight target closed-loop development seeds and changes both model fitting and execution. The anchor and source component also differ in model capacity. It therefore measures the complete selection procedure, rather than an isolated benefit of source data or the original controller alone. The city-specific MaxPressure interval still crosses zero for Los Angeles; the confirmed aggregate effect is not a separate confirmation in every city.

The external study contains two geometries. Three demand files for Jinan and additional random seeds increase coverage within those geometries; they do not increase the number of independent cities. Target-network diagnostics were used during development, although the reported final seeds were withheld. The simulator is not calibrated to field trajectories, and both adaptation and evaluation use the same physics engine. We therefore cannot infer physical deployment performance, robustness to a different simulator, or effectiveness under detector failure and unobserved demand regimes.

The word ``causal'' in the method's name refers to controlled simulator interventions and predefined input restrictions. It does not claim observational causal discovery or invariance in the physical network. Similarly, heuristic error scores do not certify safety. Incident counts in the development comparisons (Appendix~\ref{app:development-results}) demonstrate why the original zero-incident runs must not be promoted into a general safety statement.

Finally, same-protocol comparisons with transfer methods such as CrossLight, X-Light, GESA, and MetaLight remain missing. An informative next evaluation would match model capacity and target information, use the same signal executor, and include additional networks withheld from method development. It should test target-only and source-aware training with identical architectures and measure both short-horizon action quality and sustained closed-loop behavior. That design would distinguish a useful representation from a demonstrable benefit of transferred data.
